\chapter[FORMATS DE TEMPS ET DE POSITION]{Formats de temps et de
position}\label{formats-de-temps-et-de-position}

\section{Formats de temps}\label{formats-de-temps}

Les timestamps APRS s'expriment de 3 manières :

\begin{itemize}
\tightlist
\item
  Day/Hours/Minutes (DHM)
\item
  Hours/Minutes/Seconds (HMS)
\item
  Month/Day/Hours/Minutes (MDHM)
\end{itemize}

Toutes utilisent l'horloge 24 h.

\textbf{DHM} --- champ fixe de 7 caractères : 6 chiffres (jour + heure)
+ 1 indicateur (\texttt{z} ou \texttt{/}). Le groupe est
\texttt{JJHHMM}, jour du mois 01-31, heures et minutes.

\begin{itemize}
\tightlist
\item
  \texttt{092345z} = 23:45 zulu, jour 9 du mois.
\item
  \texttt{092345/} = 23:45 heure locale, jour 9.
\end{itemize}

Recommandé : n'émettre que du zulu à l'avenir. Les Status Reports
n'acceptent que le zulu.

\textbf{HMS} --- champ fixe de 7 caractères : 6 chiffres
(heures/minutes/secondes) + \texttt{h}.

\begin{itemize}
\tightlist
\item
  \texttt{234517h} = 23:45:17 zulu.
\end{itemize}

Non utilisable dans les Status Reports.

\textbf{MDHM} --- champ fixe de 8 caractères : mois (01-12), jour
(01-31), heures + minutes zulu.

\begin{itemize}
\tightlist
\item
  \texttt{10092345} = 23:45 zulu le 9 octobre.
\end{itemize}

Utilisé uniquement dans les rapports de stations météo « positionless ».

\section{Usage des timestamps}\label{usage-des-timestamps}

Quand une station émet un rapport sans timestamp, une station APRS
réceptrice peut noter en interne l'heure de réception. C'est la vision «
émission = maintenant » côté récepteur.

Quand une station émet un rapport avec timestamp, ce timestamp
représente l'heure de création côté émetteur.

Autrement dit : - rapport sans timestamp = temps réel, courant (le
récepteur note l'heure). - rapport avec timestamp = pas nécessairement
temps réel, potentiellement (très) ancien.

Quatre DTIs indiquent si le rapport contient ou non un timestamp, selon
que la station a la capacité messagerie APRS ou non :

\begin{longtable}[]{@{}lll@{}}
\toprule\noalign{}
& Sans messagerie APRS & Avec messagerie APRS \\
\midrule\noalign{}
\endhead
\bottomrule\noalign{}
\endlastfoot
Sans timestamp (courant) & \texttt{!} & \texttt{=} \\
Avec timestamp (ancien / différé) & \texttt{/} & \texttt{@} \\
\end{longtable}

Stations sans messagerie = généralement trackers autonomes ou
digipeaters. Stations sans timestamp = généralement stations fixes.

\section{Format latitude}\label{format-latitude}

Champ fixe 8 caractères : degrés + minutes décimales (2 décimales) +
\texttt{N} ou \texttt{S}.

\begin{itemize}
\tightlist
\item
  Degrés : 00 à 90.
\item
  Minutes : entier + centièmes, séparés par un point.
\end{itemize}

Ex : \texttt{4903.50N} = 49° 3' 30'' nord.

Notation générique : chaîne \texttt{ddmm.hhN}.

\section{Format longitude}\label{format-longitude}

Champ fixe 9 caractères : degrés + minutes décimales (2 décimales) +
\texttt{E} ou \texttt{W}.

\begin{itemize}
\tightlist
\item
  Degrés : 000 à 180.
\item
  Minutes : entier + centièmes.
\end{itemize}

Ex : \texttt{07201.75W} = 72° 1' 45'' ouest.

Notation générique : \texttt{dddmm.hhW}.

\section{Coordonnées de position}\label{coordonnuxe9es-de-position}

Combinaison lat + long séparées par un Symbol Table Identifier, suivies
d'un Symbol Code.

Ex : \texttt{4903.50N/07201.75W-} --- \texttt{/} = table de symboles
primaire, \texttt{-} = icône « maison ».

Voir chapitre 20 et annexe 2 pour les symboles.

\section{Ambiguïté de position}\label{ambiguuxeftuxe9-de-position}

Quand la position exacte n'est pas connue, on peut réduire la précision
en remplaçant progressivement \texttt{mm} puis \texttt{hh} de la
latitude par des espaces (\texttt{V}) :

\begin{itemize}
\tightlist
\item
  \texttt{4903.5VN} = latitude au dixième de minute.
\item
  \texttt{4903.VVN} = à la minute.
\item
  \texttt{490V.VVN} = à 10 minutes.
\item
  \texttt{49VV.VVN} = au degré.
\end{itemize}

Le niveau d'ambiguïté de la latitude s'applique automatiquement à la
longitude --- inutile d'y ajouter des \texttt{V}.

Ex : \texttt{4903.VVN/07201.75W-} = position au plus proche de la
minute. Les centièmes de minutes en lat/long peuvent prendre n'importe
quelle valeur 00-99. Bounding box : - NW : 49° 3.99' N, 72° 1.99' W - NE
: 49° 3.99' N, 72° 1.00' W - SE : 49° 3.00' N, 72° 1.00' W - SW : 49°
3.00' N, 72° 1.99' W

\section{Position nulle par
défaut}\label{position-nulle-par-duxe9faut}

Quand une station n'a pas de position spécifique (ex : Mic-E sans GPS
connecté), elle doit émettre une position nulle par défaut : 0° 0' 0''
N, 0° 0' 0'' W, avec le symbole \texttt{\textbackslash{}.} (position
inconnue).

Ex : \texttt{…0000.00N\textbackslash{}00000.00W.…}

\section{Locator Maidenhead (Grid
Square)}\label{locator-maidenhead-grid-square}

Alternative à la lat/long. Quatre manières :

\begin{itemize}
\tightlist
\item
  Dans un Status Report : \texttt{IO91SX/(} (\texttt{/-} = symbole
  maison).
\item
  Dans le Mic-E Status Text : \texttt{IO91SX/G} (\texttt{/G} = « grid
  square »).
\item
  Dans le champ Destination (obsolète) : \texttt{IO91SX}.
\item
  Dans un beacon AX.25 avec DTI \texttt{{[}} (obsolète) :
  \texttt{{[}IO91SX{]}}.
\end{itemize}

Le grid square est en forme 6 caractères ou raccourcie 4 caractères
(\texttt{IO91}).

\section{Données NMEA}\label{donnuxe9es-nmea}

APRS reconnaît les chaînes ASCII brutes conformes NMEA 0183 v2.0
provenant de GPS ou LORAN. Recommandé d'interpréter au moins :

\begin{longtable}[]{@{}ll@{}}
\toprule\noalign{}
Sentence & Signification \\
\midrule\noalign{}
\endhead
\bottomrule\noalign{}
\endlastfoot
\texttt{GGA} & Global Positioning System Fix Data \\
\texttt{GLL} & Geographic Position, Latitude/Longitude Data \\
\texttt{RMC} & Recommended Minimum Specific GPS/Transit Data \\
\texttt{VTG} & Velocity and Track Data \\
\texttt{WPT} & Way Point Location \\
\end{longtable}

\section{Altitude}\label{altitude}

Deux façons :

\begin{itemize}
\tightlist
\item
  \textbf{Dans le comment text} --- format \texttt{/A=aaaaaa},
  \texttt{aaaaaa} = altitude en \textbf{pieds}. Ex : \texttt{/A=001234}.
  Peut apparaître n'importe où dans le commentaire.
\item
  \textbf{En format Mic-E} --- champ statut Mic-E optionnel. Voir
  chapitre 10.
\end{itemize}

\begin{center}\rule{0.5\linewidth}{0.5pt}\end{center}
